\section{Introduction générale}

Ce document présente un projet académique réalisé dans le cadre de notre formation en deuxième année de licence d’informatique (L2), à l’université. Ce projet a pour vocation de mobiliser et d’approfondir un ensemble de compétences fondamentales en génie logiciel, algorithmique et programmation orientée objet. Il s’agit de concevoir et développer une simulation interactive dans laquelle des agents — plus précisément des gardiens — patrouillent une grille afin de détecter et capturer des intrus, en tenant compte de diverses contraintes comme les obstacles, la visibilité limitée .

L’objectif premier de ce projet est de simuler de manière crédible un système de surveillance automatisé, dans lequel plusieurs agents coopèrent ou s’opposent selon des règles précises. Ces agents sont dotés d’un comportement autonome, basé sur des algorithmes de planification de trajectoire , de détection d’intrus, de stratégie de poursuite ou d’évasion, et de gestion du champ de vision. Les déplacements se font dans un environnement balisé, où chaque case peut représenter un mur, une zone de passage ou un obstacle à la visibilité.

La richesse de ce projet réside dans la diversité des problématiques qu’il permet d’aborder : la gestion des interactions multi-agents, la modélisation de comportements stratégique en environnement partiellement observable, la structuration de classes Java avancées, . Chaque aspect du système a été pensé pour favoriser une compréhension plus poussée des concepts étudiés durant notre formation, tout en permettant une expérimentation visuelle concrète.

Ce projet se distingue aussi par son fort potentiel pédagogique : il met en relation directe les notions théoriques de cours (structures de données), avec des cas pratiques de développement logiciel. L’aspect visuel et interactif du jeu permet une immersion dans le fonctionnement interne du système et rend les choix d’implémentation plus parlants et motivants. Observer en temps réel les décisions d’un gardien ou la fuite d’un intrus donne un sens nouveau aux algorithmes que nous implémentons.

Au fil de ce rapport, nous détaillerons la méthodologie suivie, les choix d’architecture logicielle, les algorithmes implémentés, les stratégies développées par les agents, des pistes d’optimisation, et des perspectives d’amélioration pour des versions futures du projet. Ce travail représente non seulement une réalisation technique complète, mais aussi une démarche de réflexion sur les méthodes de génie logiciel adaptées aux systèmes complexes.
